%!TEX root = ../../main.tex
%% ===============================================================
%% 9.3 조건부 예측 가능성
%% 파일: chapters/09_results/03_conditional.tex
%% 상위: chapters/09_results/chapter.tex
%% 정의 라벨: sec:results-conditional
%% 참조 라벨: tab:conditional
%% 포함 객체: tables/tab_conditional
%% 인용: -
%% 상태: 초안 (docs/OUTLINE.md 참고)
%% ===============================================================

\section{조건부 예측 가능성}
\label{sec:results-conditional}

%% ---------------------------------------------------------------
%% table 객체: tab:conditional — 하위 집단별 LightGBM 시험 AUC
%% 파일: tables/tab_conditional.tex   사용처: chapters/09_results/03_conditional.tex
%% ---------------------------------------------------------------
\begin{table}[htbp]
  \centering
  \caption{하위 집단별 LightGBM 시험 AUC}
  \label{tab:conditional}
  \small
  \begin{tabular}{@{}llc@{}}
    \toprule
    축 & 하위 집단 & AUC \\
    \midrule
    \multirow{3}{*}{게임 국면} & 초반 ($<15$ 분) & \aucEarly{} \\
     & 중반 (15--25 분) & \aucMid{} \\
     & 후반 ($>25$ 분) & \aucLate{} \\
    \midrule
    \multirow{3}{*}{온셋 골드 격차} & 대등 ($|\Delta| < 2{,}000$) & \aucClose{} \\
     & 보통 & \aucModerate{} \\
     & 일방 ($|\Delta| > 5{,}000$) & \aucOneSided{} \\
    \bottomrule
  \end{tabular}
\end{table}


표 \ref{tab:conditional}은 RQ4에 답한다. 예측 가능성은 게임 국면에 따라 $\aucEarly{} \to \aucMid{} \to \aucLate{}$로, 골드 격차에 따라 $\aucClose{} \to \aucModerate{} \to \aucOneSided{}$으로 단조 증가한다. 후반과 일방적 교전에서는 누적된 자원 격차가 프레임 수준 상태에 이미 반영되어 있어 분 단위 관측으로도 결과가 거의 결정된다. 반대로 대등한 교전에서는 여러 신경망 표현이 우연 수준($.5$)에 근접한다. 대등한 교전의 결과를 가르는 것은 스킬 적중, 진입 각도, 쿨다운 상태처럼 초 단위 이하에서 일어나는 사건인데, 이는 공개 텔레메트리에 존재하지 않는다. 즉 관측된 AUC 상한은 모델의 한계가 아니라 \emph{정보 입도}의 한계로 해석된다.
